% Replace this subsection with the completed study, not merely a populated table.
% Keep observed field facts separate from model outputs and synthetic results.
\subsection{Erlangen, Germany: local evaluation to be completed}
\label{sec:erlangen}
\AuthorQuery{Insert the actual Erlangen study design, results and interpretation
here. No local field measurements or city-scale performance results were
provided for this draft. The following is an evaluation plan and reporting
scaffold, not a claim that these procedures have already been performed.}

The local evaluation will examine route comparisons on a fixed OSM extract
covering the selected study area in Erlangen, Germany. The report should
identify the extraction date, geographic boundary, input checksum, graph size
and any preprocessing. Origin--destination pairs should be specified using
an auditable sampling rule, rather than choosing only favourable examples.
Both successful and failed requests, including snapping failures, should be
reported. The approximately 2~km case motivating the output redesign can be
included, but its actual route, observations and timing must be supplied.

For each successful request, compare the selected route with both the
no-preference standard and the distance-only route under matched hard
constraints. Keep endpoints, input snapshot and assessment profile fixed.
Report distance, grouped feature encounters, known preference conflicts,
unknown observations, coverage and feasibility; score differences are defined
only for feasible routes under the same measuring rule. Distinguish
model-reported obstacles from independently audited physical features, using
a stated correspondence rule between audited features and OSM elements.
Audit dates, measurement methods, accessible alternatives and uncertainty
should accompany any assertion about physical suitability. Do not ask a
participant to traverse an alternative that violates their requirements merely
to complete a comparison.

Runtime evaluation should separate loading, graph construction, search,
diagnostics, figures and archive creation, and record hardware, software
versions, cold and warm conditions, repetitions and memory-measurement methods.
Summaries should retain the paired structure, state their denominators and
show variability. If repeated requests or several profiles share an
origin--destination pair, uncertainty estimates must not treat those
observations as independent without justification.

\begin{table}[!ht]
\small
\begin{tabularx}{\textwidth}{@{}>{\raggedright\arraybackslash}p{0.43\textwidth}Y@{}}
\toprule
\textbf{Required Erlangen information} & \textbf{Author-supplied result} \\
\midrule
Study boundary; extract date/hash; nodes/edges & \Pending \\
Route-pair selection; requested/successful/failed counts & \Pending \\
Profiles, hard requirements and cost parameters & \Pending \\
Independent audit method, dates and feature correspondence & \Pending \\
Paired distances, encounters, conflicts and unknowns & \Pending \\
Feasibility and common-profile scores/coverage & \Pending \\
Runtime/memory, hardware, repetitions and variability & \Pending \\
Field agreement, uncertainty and unfavourable cases & \Pending \\
Data availability, OSM attribution and privacy safeguards & \Pending \\
\bottomrule
\end{tabularx}
\caption{Placeholder for the Erlangen evaluation. ``Pending'' means no result
has been supplied; it does not mean zero, no obstacle, or a successful test.}
\label{tab:erlangen}
\end{table}

\AuthorQuery{Replace the plan with what was actually done and report the
results, including null or adverse findings. Add an appropriately attributed
Erlangen map only after choosing the real study area. If participants,
identifiable locations or personal mobility data were involved, state the
applicable ethics approval or exemption, consent and privacy arrangements;
do not claim approval or exemption without verification. Update the abstract,
impact and conclusions to reflect only the supplied evidence.}